% ---
% title: Sharing Work Online 
% date: 2026-06
% category:
% nodes: reflections, reading
% links:
% ---
\documentclass{article}
\usepackage[utf8]{inputenc}
\usepackage{amsmath, amssymb, amsthm}
\usepackage{geometry}
\usepackage{graphicx}
\usepackage{hyperref}

\geometry{a4paper, margin=1in}

\title{Sharing Work Online}
\date{June 2026}
\author{Daksh Mehta}

\begin{document}


I recently read Show Your Work! by \href{https://austinkleon.com}{Austin Kleon}. I find it difficult to share things online and often wonder what the point is. Messy, unpolished things should remain private and only polished final outputs should make it to the LinkedIn profile or wherever you choose to share things about yourself. The book encourages readers to share their work more regularly. The author talks about how digital spaces have removed traditional mediums and allow people to connect directly to an audience – and creators are beginning to discover that progress and regular documentation is also valuable. Not only to the audience but also to yourself. The book sells you on the idea that regular documentation of your process is a beneficial thing to do. There are selfish reasons (a record frames you as a consistent and introspective worker which can be appealing for employers), spiritual reasons (the art of reflecting regularly and aligning yourself with your authentic self) and a reason for those who don't like reasons (why not? what is there to lose?). 

I find that forcing yourself to sit down and write about something you've done – whether you do end up putting it anywhere or not – organises thoughts in a very neat way, giving a cross-sectional slice of who you were at that period of time. Even when I'm writing about work, I can look back and remember all the little challenges and wins and how they made me feel. More than anything, I’d like to make writing a regular habit. I’m not entirely sure what getting better at writing means in practice, but I think consistency and a growing collection of things to look back on will help me figure that out.

I am also trying to minimise my connectivity to the online world by being more intentional with my usage, albeit rather painfully slowly. I think paradoxically, creating and publishing to online spaces with the eventual goal of creating on the same order as consuming will help claw back some control. 

The later chapters in the book talk about growing your presence, receiving feedback and building meaningful connections. I might revisit this when I get further along my journey. I am treating this as an experiment, perhaps nothing comes of it. But if I write regularly, think a little more carefully and leave behind a record of what I cared about at different points, that seems reason enough to continue. 


\end{document}
